\section{Планарная граница и продуктовое исчисление ширин}
\label{sec:product}

Оценки в размерностях $9$ и $10$ получаются без единого численного шага: из
одного точного факта о блоке $E_8/2401$ и одного правила сложения ширин.

\subsection{Планарная граница для эйзенштейновых конструкций}

Решётка $\Lambda\subset\R^n$ \emph{эйзенштейнова} (имеет
$\Z[\omega]$-структуру, $\omega=e^{2\pi i/3}$), если на ней задано умножение
на $\omega$, действующее изометрией порядка~$3$ без неподвижных ненулевых
векторов; тогда $\Lambda$ --- $\Z[\omega]$-модуль, и для
$\alpha\in\Z[\omega]$ подрешётка $\alpha\Lambda$ имеет индекс $|\alpha|^n$.
Конструкции АБПР $7^{n/2}$ цветов в размерностях $n=4,6,8,24$ --- это
$\Gamma=(3+\omega)\Lambda$ ($|3+\omega|^2=7$) на $D_4$, $E_6^*$, $E_8$ и
решётке Лича~\cite{ABPR}.

\begin{proposition}[планарная нижняя граница]\label{prop:planar}
Пусть $\Lambda\subset\R^n$ --- эйзенштейнова решётка. Тогда для каждого
ненулевого $w\in\Lambda$
\[
D\bigl((3+\omega)w\bigr)\ \ge\ \sqrt{7/3}\,\|w\|,
\qquad\text{в частности}\qquad
D\bigl((3+\omega)\Lambda\bigr)\ \ge\ \sqrt{7/3}\,\lambda_1(\Lambda),
\]
где $\lambda_1$ --- длина кратчайшего ненулевого вектора.
\end{proposition}

\begin{proof}
Векторы $\pm w,\ \pm\omega w,\ \pm(1+\omega)w$ лежат в $\Lambda$ и имеют одну
норму $\|w\|$ (умножение на обратимые элементы $\Z[\omega]$ --- изометрия
плоскости $P=\mathrm{span}(w,\omega w)$). Соответствующие шесть полупространств
$\langle x,u\rangle\le\|u\|^2/2$ содержат $V_0$, а их нормали лежат в $P$,
поэтому ортопроекция $\pi_P(V_0)$ содержится в правильном шестиугольнике $H$ ---
ячейке Вороного решётки $\Z[\omega]w\cong A_2$ масштаба $\|w\|$ внутри $P$.
Точка $v/2=(3+\omega)w/2$ лежит в $P$, проекция $1$-липшицева, и
\[
\dist\bigl(v/2,\,V_0\bigr)\ \ge\ \dist_P\bigl(v/2,\,H\bigr)
=\sqrt{7/12}\,\|w\|:
\]
в координатах плоскости при $\|w\|=1$ точка $(3+\omega)/2=(5/4,\sqrt3/4)$
лежит в нормальном конусе вершины $(1/2,\,1/(2\sqrt3))$ шестиугольника с
фасетами на расстоянии $1/2$, и квадрат расстояния до неё равен
$9/16+1/48=7/12$. По лемме~\ref{lem:mid}
$D(v)=2\dist(v/2,V_0)\ge\sqrt{7/3}\,\|w\|$; всякий ненулевой вектор
$(3+\omega)\Lambda$ имеет вид $(3+\omega)w$ с $w\ne0$, и $\|w\|\ge\lambda_1$.
\end{proof}

Для блока $E_8/2401$ этого достаточно. Решётка $E_8$ эйзенштейнова \cite{ConwaySloane,ABPR}; при минимальной
норме $\lambda_1^2=2$ её радиус покрытия равен~$1$ \cite{ConwaySloane}, так
что $\diam V_0=2$.
В нормировке $\lambda_1^2=3$ получаем $\diam^2V_0=6$, а по
предложению~\ref{prop:planar} $D\bigl((3+\omega)E_8\bigr)\ge\sqrt7$. Итак,
\begin{equation}\label{eq:e8width}
d\bigl(E_8,(3+\omega)E_8\bigr)^2\ \ge\ \tfrac76,\qquad
\bigl[E_8:(3+\omega)E_8\bigr]=7^4=2401 .
\end{equation}
(На самом деле здесь равенство --- минимум достигается на образах минимальных
векторов, см.~\cite{Widths}; для оценок ниже нужна только нижняя граница.)

\subsection{Правило сложения ширин}

\begin{proposition}[продуктовое исчисление]\label{prop:product}
Пусть $\Lambda=\bigoplus_{i=1}^{s}\Lambda_i$ --- ортогональная прямая сумма
решёток $\Lambda_i\subset\R^{n_i}$, и $\Gamma=\bigoplus_i\Gamma_i$ с
$\Gamma_i\subseteq\Lambda_i$. Обозначим $d_i=D(\Gamma_i)/\diam V_0(\Lambda_i)$
и $k_i=[\Lambda_i:\Gamma_i]$. Тогда после независимого масштабирования блоков
раскраска по $\Gamma$ правильна для интервала $[1,\ell]$ при всех
\[
\ell<\Bigl(\sum_i 1/d_i^2\Bigr)^{-1/2};
\]
в частности, $\chi\bigl(\R^{n},[1,\ell]\bigr)\le\prod_i k_i$, $n=\sum_i n_i$,
как только $\sum_i 1/d_i^2\le1$. Оптимум достигается при
$\diam^2 V_0(\Lambda_i)\propto 1/d_i^2$.
\end{proposition}

\begin{proof}
Ячейка Вороного прямой суммы есть произведение ячеек, поэтому квадрат
диаметра $V_0(\Lambda)$ равен $\sum_i\diam^2 V_0(\Lambda_i)$. Далее, для
$v=(v_1,\dots,v_s)$
\[
\dist(v/2,V_0)^2=\sum_i\dist\bigl(v_i/2,V_0(\Lambda_i)\bigr)^2,
\qquad\text{то есть}\qquad D(v)^2=\sum_i D_i(v_i)^2 .
\]
Так как $D_i(0)=0$, минимум по
$\Gamma\setminus\{0\}$ достигается на векторе с единственной ненулевой
компонентой, и условие $D(\Gamma)\ge\ell\diam V_0$ распадается на
$s$ условий
\[
d_i^2\,s_i\ \ge\ \ell^2\sum_j s_j,\qquad s_i:=\diam^2 V_0(\Lambda_i).
\]
Блоки можно масштабировать независимо, поэтому положим $s_i=S/d_i^2$; тогда
каждое условие обращается в равенство при $\ell^2=S/\sum_j s_j
=\bigl(\sum_j 1/d_j^2\bigr)^{-1}$, и по предложению~\ref{prop:crit} раскраска
правильна при всех меньших $\ell$. Любой другой выбор $s_i$ даёт не большее
$\ell$: минимум по $i$ величины $d_i^2 s_i/\sum_j s_j$ максимизируется именно
при пропорциональности. Наконец, правая часть монотонна по каждому $d_i$,
поэтому вместо точных ширин достаточно нижних оценок.
\end{proof}

Ширина оказывается не отчётом о качестве раскраски, а \emph{расходуемым
ресурсом}: блок с шириной $d_i$ стоит $1/d_i^2$ единиц бюджета, равного~$1$.
Существенно, что блоки могут иметь \emph{разный индекс}: именно смесь дорогого
широкого блока с дешёвым узким и работает.

\subsection{Размерности $10$ и $9$}

\begin{keybox}
\begin{theorem}\label{thm:dim10}
$\displaystyle\chi\Bigl(\R^{10},\bigl[1,\ell\bigr]\Bigr)\ \le\ 2401\cdot19=45619$
при всех $\ell<\sqrt{217/214}$; в частности, $\chi(\R^{10})\le45619$.
\end{theorem}
\end{keybox}

\begin{proof}
Берём два блока.

\emph{Блок $E_8$.} $\Gamma_1=(3+\omega)E_8$, индекс $2401$; по~\eqref{eq:e8width}
$d_1^2\ge7/6$, цена блока $1/d_1^2\le6/7$.

\emph{Плоский блок.} $\Lambda_2=aA_2$ --- шестиугольная решётка,
$\Gamma_2=(5+2\omega)\,aA_2$, индекс $N(5+2\omega)=19$. В плоскости ячейка
Вороного \emph{и есть} шестиугольник, поэтому $D$ вычисляется точно: при
$a=1$ ближайшая к точке $(5+2\omega)/2=(2,\sqrt3/2)$ точка шестиугольника с
фасетами на расстоянии $1/2$ --- его вершина $(1/2,1/(2\sqrt3))$, квадрат
расстояния до неё равен $9/4+1/3=31/12$, а по всем векторам $(5+2\omega)w$,
$w\in A_2\setminus\{0\}$, минимум достигается при $\|w\|=1$. Отсюда
$D(\Gamma_2)=2a\sqrt{31/12}$, $\diam V_0=2a/\sqrt3$ и $d_2^2=31/4$, цена
$1/d_2^2=4/31$.

Сумма цен $6/7+4/31=214/217<1$, и предложение~\ref{prop:product} даёт
$\ell<\sqrt{217/214}$ при $\prod k_i=2401\cdot19=45619$.
\end{proof}

Прежняя оценка в этой размерности --- $3^{10}=59049$ \cite{ABPR}, так что
выигрыш составляет $22{,}7\%$, и $45619$ --- первая оценка в $\R^{10}$ ниже
классической башни $3^n$. Оптимальный масштаб плоского блока
$a=\sqrt7\big/\bigl(2\sqrt{31/12}\bigr)=0{,}8230549\ldots$; сквозная численная
проверка построенной десятимерной решётки (индекс ровно $45619$,
$d=1{,}006984951\ldots$, совпадение с $\sqrt{217/214}$ до $10^{-9}$) в
доказательстве не участвует и служит лишь контролем.

\begin{keybox}
\begin{corollary}\label{cor:dim9}
$\displaystyle\chi\Bigl(\R^9,\bigl[1,\ell\bigr]\Bigr)\ \le\ 2401\cdot4=9604$
при всех $\ell<\sqrt{63/61}$; в частности, $\chi(\R^9)\le9604$.
\end{corollary}
\end{keybox}

\begin{proof}
Первый блок --- тот же $E_8/2401$ с ценой не выше $6/7$. Второй --- одномерный:
$\Lambda_2=\Z$, $\Gamma_2=4\Z$, четыре цвета. Ячейка Вороного --- отрезок
длины~$1$, две ближайшие одноцветные ячейки суть $[-1/2,1/2]$ и $[7/2,9/2]$,
расстояние между ними равно~$3$, поэтому $d_2=3$ и цена $1/d_2^2=1/9$. Сумма
цен $6/7+1/9=61/63<1$, и предложение~\ref{prop:product} даёт
$\ell<\sqrt{63/61}=1{,}016261\ldots$ при $\prod k_i=2401\cdot4=9604$.
\end{proof}

Три цвета не проходят: там $d_2=2$, цена $1/4$, и $6/7+1/4=31/28>1$ ---
бюджет $1/7$, оставленный блоком $E_8/2401$, требует ровно $d_2\ge\sqrt7$.
Оценка улучшает $17253$ работы~\cite{ABPR} на $44\,\%$. Тот же счёт с блоком
$\Lambda_{24}/7^{12}$ (у решётки Лича также $\lambda_1^2=4$ при радиусе
покрытия $\sqrt2$, то есть та же ширина $\sqrt{7/6}$) даёт
$\chi(\R^{25})\le4\cdot7^{12}$ и $\chi(\R^{26})\le19\cdot7^{12}$ против
$3^{25}$ и $3^{26}$; подробности, а также численные конструкции с меньшим
числом цветов в $\R^9$ и $\R^{10}$, чьи сертификаты диаметра пока выполнены в
плавающей точке, --- в~\cite{Widths}.
